\documentclass{article}
\usepackage{amsmath}
\usepackage{geometry}
\usepackage{xcolor}
\geometry{letterpaper, margin=1in}
\title{EMCH 290 Quiz\#1 Solution Manual}
\date{}

\begin{document}

\maketitle

\section*{Question 1 [10 points]}

\textbf{Question:} A satellite weighs 9800 N on Earth where \( g = 9.8 \, \text{m/s}^2 \). What is the weight of the satellite, in N, in space when the acceleration of gravity is \( 0.25 \, \text{m/s}^2 \)?
\newline \newline
\textbf{Solution:} 
The weight \( W \) of an object is given by \( W = mg \), where \( m \) is the mass and \( g \) is the gravitational acceleration. The mass of the satellite can be calculated from its weight on Earth as follows:

\[
m = \frac{W_{\text{Earth}}}{g_{\text{Earth}}} = \frac{9800}{9.8} = 1000 \, \text{kg}
\]
\newline
The weight \( W_{\text{Space}} \) of the satellite in space is given by:

\[
W_{\text{Space}} = m \times g_{\text{Space}} = 1000 \times 0.25 = 250 \, \text{N}
\]
\newline
\textbf{Points Breakdown:}
\begin{enumerate}
    \renewcommand{\labelenumi}{\alph{enumi})}
    \item Correctly calculating mass: 3 points
    \item Correctly applied formula: 3 points
    \item Correct final answer: 4 points
\end{enumerate}



\newpage
\section*{Question 2 [20 points]}

\textbf{Question:} A bottle made of 12 kg steel contains 2 kmole of liquid propane in it. It is on a truck, and a strap breaks. The bottle accelerates horizontally with an acceleration of \(5 \, \text{m/s}^2\). What is the needed force for this acceleration? (Molecular weight of propane is \( M = 44.1 \, \text{kg/kmol} \))
\newline \newline
\textbf{Solution:} 
The total mass \(m_{\text{Total}}\) is the mass of the steel bottle plus the mass of the propane:

\[
m_{\text{Total}} = m_{\text{Steel}} + m_{\text{Propane}} = 12 \, \text{kg} + (2 \, \text{kmol} \times 44.1 \, \text{kg/kmol}) = 12 \, \text{kg} + 88.2 \, \text{kg} = 100.2 \, \text{kg}
\]
\newline
The force \( F \) needed for this acceleration is \( F = m_{\text{Total}} \times a \):

\[
F = 100.2 \, \text{kg} \times 5 \, \text{m/s}^2 = 501 \, \text{N}
\]
\newline
\textbf{Points Breakdown:}
\begin{enumerate}
    \renewcommand{\labelenumi}{\alph{enumi})}
    \item Calculating the mass of propane: 5 points
    \item Calculating total mass: 5 points
    \item Applying \( F = ma \): 5 points
    \item Correct answer: 5 points
\end{enumerate}



\newpage
\section*{Question 3 [15 points]}
\textbf{Question:} You dive 10 m deep in a pool. What is the absolute pressure there?
\newline \newline
\textbf{Solution:} 
The pressure difference due to a fluid column is given by \( \Delta P = \rho g h \), where \( \rho \) is the fluid density, \( g \) is the gravitational acceleration, and \( h \) is the height of the fluid column.
\newline \newline
For water, \( \rho = 1000 \, \text{kg/m}^3 \) and \( g = 9.8 \, \text{m/s}^2 \). At 10 m deep, the pressure difference \( \Delta P \) is:
\[
\Delta P = 1000 \times 9.8 \times 10 = 98000 \, \text{Pa} = 98 \, \text{kPa}
\]
The absolute pressure is the atmospheric pressure plus this difference: \( P_{\text{abs}} = P_{\text{atm}} + \Delta P = 101 \, \text{kPa} + 98 \, \text{kPa} = 199 \, \text{kPa} \).
\newline \newline
\textbf{Points Breakdown:}
\begin{enumerate}
    \renewcommand{\labelenumi}{\alph{enumi})}
    \item Correctly applying the formula for pressure difference: 5 points
    \item Correct calculation of \( \Delta P \): 5 points
    \item Correct answer for \( P_{\text{abs}} \): 5 points
\end{enumerate}




\newpage
\section*{Question 4 [20 points]}
\textbf{Question:} What is the absolute pressure (kPa) at the bottom of a 5 m tall column of fluid if the fluid is a) water b) glycerin c) gasoline?
\newline \newline
\textbf{Solution:} 
\newline
To find the absolute pressure at the bottom of each column, we use the formula \( \Delta P = \rho g h \) and then add atmospheric pressure \( P_{\text{atm}} = 101 \, \text{kPa} \) to get \( P_{\text{abs}} \). 
\newline We assume the container is open to the atmosphere.
\newline \newline
\textbf{Step 1: Water}
\newline
\( \Delta P = 1000 \, \text{kg/m}^3 \times 9.8 \, \text{m/s}^2 \times 5 \, \text{m} = 49,000 \, \text{Pa} = 49 \, \text{kPa} \)
\newline
\( P_{\text{abs}} = 49 \, \text{kPa} + 101 \, \text{kPa} = 150 \, \text{kPa} \)
\newline \newline
\textbf{Step 2: Glycerin}
\newline
\( \Delta P = 1260 \, \text{kg/m}^3 \times 9.8 \, \text{m/s}^2 \times 5 \, \text{m} = 61,740 \, \text{Pa} = 61.74 \, \text{kPa} \)
\newline
\( P_{\text{abs}} = 61.74 \, \text{kPa} + 101 \, \text{kPa} = 162.74 \, \text{kPa} \)
\newline \newline
\textbf{Step 3: Gasoline}
\newline
\( \Delta P = 750 \, \text{kg/m}^3 \times 9.8 \, \text{m/s}^2 \times 5 \, \text{m} = 36,750 \, \text{Pa} = 36.75 \, \text{kPa} \)
\newline
\( P_{\text{abs}} = 36.75 \, \text{kPa} + 101 \, \text{kPa} = 137.75 \, \text{kPa} \)
\newline \newline
\textbf{Points Breakdown:}
\begin{enumerate}
    \renewcommand{\labelenumi}{\alph{enumi})}
    \item Identification and correct application of the formula \( \Delta P = \rho g h \): 5 points
    \item Accurate calculation of \( \Delta P \) for water: 2 points
    \item Correct \( P_{\text{abs}} \) for water: 2 points
    \item Accurate calculation of \( \Delta P \) for glycerin: 2 points
    \item Correct \( P_{\text{abs}} \) for glycerin: 2 points
    \item Accurate calculation of \( \Delta P \) for gasoline: 2 points
    \item Correct \( P_{\text{abs}} \) for gasoline: 2 points
    \item Overall coherence and step-by-step flow of the solution: 3 points
\end{enumerate}





\newpage
\section*{Question 5 [20 points]}
\textbf{Question:} A piston with additional weight has been suspended on top of a cylinder containing a gas. The weight of the piston and weights is a combined 100 kg. Find the work that is done on the system in joule if the change in piston location is 0.51m.
\newline \newline
\textbf{Solution:}
\newline
To find the work done, \( W \), we use the formula \( W = Fd \), where \( F \) is the force and \( d \) is the distance.
\newline \newline
\textbf{Step 1:} Calculation of Force \( F \) 
\newline
First, we calculate the force \( F \) due to gravity acting on the piston and weights. \( F = mg = 100 \times 9.8 = 980 \, \text{N} \).
\newline \newline
\textbf{Step 2:} Application of Work Formula
\newline
Next, we plug this force \( F \) into the work formula \( W = Fd \). Here, the distance \( d \) is 0.51m.
\newline  \newline
\textbf{Step 3:} Final Calculation
\newline
Finally, we find \( W = 980 \times 0.51 = 499.8 \, \text{J} \).
\newline \newline
\textbf{Points Breakdown:}
\begin{enumerate}
    \renewcommand{\labelenumi}{\alph{enumi})}
    \item Identification of the correct formula \( W = Fd \) to use for work calculation: 2 points
    \item Accurate calculation of the force \( F = mg \): 4 points
    \item Correct substitution of \( F \) into the work formula \( W = Fd \): 3 points
    \item Correct multiplication to find \( W \): 3 points
    \item Final answer with correct units (Joules): 3 points
    \item Overall coherence and flow of the solution: 5 points
\end{enumerate}




\newpage
\section*{Question 6 [15 points]}
\textbf{Question:} State whether each of the following systems is in Thermal equilibrium, Chemical equilibrium, or Mechanical equilibrium.
\newline \newline
\textbf{Solution:}
\newline
a) Snow resting on the roof of your house on a very cold day: Mechanical \& Chemical equilibrium
\newline
Explanation: Here, the snow isn't melting or sublimating, so it's in chemical equilibrium. Also, it's stationary and not subject to unbalanced forces, indicating mechanical equilibrium.
\newline \newline
b) A bicycle rolling along a level road: Mechanical Equilibrium
\newline
Explanation: A bicycle rolling along a level road at a constant speed is in mechanical equilibrium. This means that the forces acting on the bicycle are balanced, resulting in no net force to accelerate or decelerate it. The main forces involved include gravity, which is balanced by the road's normal force, and friction, which is countered by the bicycle's propulsive force when pedaling or by its existing motion if not actively pedaling. This equilibrium allows the bicycle to maintain its constant speed.
\newline \newline
c) A cup of coffee sitting on your desk: Chemical \& Mechanical Equilibrium
\newline
Explanation: The coffee, if undisturbed, is not undergoing any chemical changes; hence, it's in chemical equilibrium. It's also not moving, so it's in mechanical equilibrium. However, if it's cooling down, it's not in thermal equilibrium.
\newline \newline
\textcolor{red}{
*Note that multiple interpretations are acceptable, provided that sufficient assumptions or evidence is shown.
}
\newline \newline
\textbf{Points Breakdown:}
\begin{enumerate}
    \renewcommand{\labelenumi}{\alph{enumi})}
    \item Correct identification of equilibrium type for snow and valid explanation: 5 points
    \item Correct identification of equilibrium type for bicycle and valid explanation: 5 points
    \item Correct identification of equilibrium type for coffee and valid explanation: 5 points
\end{enumerate}


\end{document}


\end{document}
